% Generated by the blueprint tool (blueprint-pdf/1) from a blueprint snapshot; edit the
% Markdown sources, not this file.  Compile with `latexmk -pdf` or by running
% `pdflatex` twice (the second run settles the contents and the page references).
% Images, if any, are read from assets/ next to this file.
\documentclass[11pt]{article}
\usepackage[T1]{fontenc}
\usepackage[utf8]{inputenc}
\usepackage{lmodern}
\usepackage{textcomp}
\usepackage[a4paper,margin=2.5cm,headheight=14pt]{geometry}
\usepackage{amsmath,amssymb,amsthm}
\usepackage{mathtools}
\usepackage{bm}
\usepackage{graphicx}
\usepackage{xcolor}
\usepackage{array,booktabs,tabularx}
\usepackage{fancyvrb}
\usepackage[normalem]{ulem}
\usepackage{microtype}
% Optional, and only with their fonts: a small TeX installation has the
% packages without the fonts, which fails at the first formula.
\IfFileExists{mathrsfs.sty}{\IfFileExists{rsfs10.tfm}{\usepackage{mathrsfs}}{}}{}
\IfFileExists{stmaryrd.sty}{\IfFileExists{stmary10.tfm}{\usepackage{stmaryrd}}{}}{}
\IfFileExists{cancel.sty}{\usepackage{cancel}}{}
\usepackage[hyphens]{url}
\usepackage[pdfusetitle,bookmarksnumbered=false]{hyperref}
\usepackage{bookmark}
\definecolor{bplink}{HTML}{1F4E9C}
\definecolor{bpmuted}{HTML}{5F676E}
\hypersetup{colorlinks=true,linkcolor=bplink,urlcolor=bplink,citecolor=bplink}
\setlength{\parindent}{0pt}
\setlength{\parskip}{0.55em plus 0.2em minus 0.1em}
% Lists as tight as the page's.
\makeatletter
\def\@listi{\leftmargin\leftmargini\parsep 0pt\topsep 0.3em plus 0.1em\itemsep 0.1em plus 0.1em}
\let\@listI\@listi
\def\@listii{\leftmargin\leftmarginii\labelwidth\leftmarginii\advance\labelwidth-\labelsep\parsep 0pt\topsep 0.1em\itemsep 0.1em}
% An ordered list that does not start at 1.
\newcommand{\bpstartat}[1]{\setcounter{\@enumctr}{\numexpr#1-1\relax}}
\makeatother
\setcounter{tocdepth}{3}
\emergencystretch=3em
\allowdisplaybreaks

% A numbered statement, set as a paper sets one: "Theorem 2.1 (Title)."
% #1 is the lead ("Theorem~2.1"), #2 what follows it in roman.
\newenvironment{bpstatement}[2]{\par\addvspace{0.6\baselineskip}\noindent\textbf{#1}#2\textbf{.}\enspace\ignorespaces}
  {\par\addvspace{0.4\baselineskip}}
\newenvironment{bpstatementit}[2]{\begin{bpstatement}{#1}{#2}\itshape\ignorespaces}{\end{bpstatement}}
% A step: a relation with prose of its own, read under the object it starts from.
\newenvironment{bpstep}[1]{\par\addvspace{0.3\baselineskip}\begin{list}{}{\leftmargin=1.6em\rightmargin=0pt\topsep=0pt}\item\relax
  {\small\color{bpmuted}#1}\par\nopagebreak}{\end{list}\par}
% Notes about where an object is written out.
\newenvironment{bpnote}{\par\small\color{bpmuted}\itshape}{\par}
\newcommand{\bpomitted}[1]{\par{\small\color{bpmuted}\itshape #1\par}}
\newcommand{\bpbodyheading}[1]{\par\addvspace{0.4\baselineskip}\noindent\textbf{#1}\par\nopagebreak}
\newcommand{\bprule}{\par\medskip\noindent\hfil\rule{0.3\linewidth}{0.4pt}\hfil\par\medskip}
\newcommand{\bpboundary}[1]{\par\nopagebreak{\small\color{bpmuted}#1\par}\nopagebreak}
% A reference to something this export leaves out, and one that resolves nowhere.
\newcommand{\bpout}[1]{#1\textsuperscript{\hyperlink{bp-outside}{\ensuremath{\ast}}}}
\newcommand{\bpbroken}[1]{\textcolor{red!70!black}{[#1?]}}
\newcommand{\bpslug}[1]{\texttt{#1}}
% A character the export knows no way to set.
\newcommand{\bpglyph}[1]{\fbox{\scriptsize #1}}
% Images, at their own size unless they are wider than the line.
\newsavebox{\bpimagebox}
\newcommand{\bpimage}[1]{\sbox{\bpimagebox}{\includegraphics{#1}}%
  \ifdim\wd\bpimagebox>\linewidth\resizebox{\linewidth}{!}{\usebox{\bpimagebox}}\else\usebox{\bpimagebox}\fi}
\newcommand{\bpmissingimage}[1]{\fbox{\parbox{0.8\linewidth}{\centering\small\color{bpmuted}Image not included: #1}}}
\DefineVerbatimEnvironment{bpverbatim}{Verbatim}{fontsize=\small,xleftmargin=1em}
\renewcommand{\qedsymbol}{\ensuremath{\blacksquare}}
\hypersetup{pdftitle={Compactness through finite intersections},pdfsubject={The whole blueprint; all levels},pdfcreator={blueprint}}
% KaTeX commands the website knows
\providecommand{\R}{\mathbb{R}}\providecommand{\N}{\mathbb{N}}\providecommand{\Z}{\mathbb{Z}}
\providecommand{\Q}{\mathbb{Q}}\providecommand{\C}{\mathbb{C}}
\providecommand{\Reals}{\mathbb{R}}\providecommand{\reals}{\mathbb{R}}\providecommand{\Complex}{\mathbb{C}}
\providecommand{\cnums}{\mathbb{C}}\providecommand{\natnums}{\mathbb{N}}\providecommand{\Bbb}[1]{\mathbb{#1}}
\providecommand{\bold}[1]{\mathbf{#1}}\providecommand{\frak}[1]{\mathfrak{#1}}\providecommand{\mathbbm}[1]{\mathbb{#1}}
\providecommand{\lang}{\langle}\providecommand{\rang}{\rangle}\providecommand{\bra}[1]{\langle #1\rvert}
\providecommand{\ket}[1]{\lvert #1\rangle}\providecommand{\braket}[1]{\langle #1\rangle}
\providecommand{\set}[1]{\{#1\}}\providecommand{\Set}[1]{\left\{#1\right\}}
\providecommand{\infin}{\infty}\providecommand{\isin}{\in}\providecommand{\plusmn}{\pm}\providecommand{\sub}{\subset}
\providecommand{\sube}{\subseteq}\providecommand{\supe}{\supseteq}\providecommand{\exist}{\exists}
\providecommand{\bull}{\bullet}\providecommand{\larr}{\leftarrow}\providecommand{\rarr}{\rightarrow}
\providecommand{\lrarr}{\leftrightarrow}\providecommand{\Larr}{\Leftarrow}\providecommand{\Rarr}{\Rightarrow}
\providecommand{\Lrarr}{\Leftrightarrow}\providecommand{\uarr}{\uparrow}\providecommand{\darr}{\downarrow}
\providecommand{\harr}{\leftrightarrow}\providecommand{\lArr}{\Leftarrow}\providecommand{\rArr}{\Rightarrow}
\providecommand{\hArr}{\Leftrightarrow}\providecommand{\alef}{\aleph}\providecommand{\alefsym}{\aleph}
\providecommand{\real}{\Re}\providecommand{\image}{\Im}\providecommand{\weierp}{\wp}
\providecommand{\Alpha}{\mathrm{A}}\providecommand{\Beta}{\mathrm{B}}\providecommand{\Epsilon}{\mathrm{E}}
\providecommand{\Zeta}{\mathrm{Z}}\providecommand{\Eta}{\mathrm{H}}\providecommand{\Iota}{\mathrm{I}}
\providecommand{\Kappa}{\mathrm{K}}\providecommand{\Mu}{\mathrm{M}}\providecommand{\Nu}{\mathrm{N}}
\providecommand{\Omicron}{\mathrm{O}}\providecommand{\Rho}{\mathrm{P}}\providecommand{\Tau}{\mathrm{T}}
\providecommand{\Chi}{\mathrm{X}}\providecommand{\omicron}{o}\providecommand{\gt}{>}\providecommand{\lt}{<}
\providecommand{\And}{\&}\providecommand{\hdashline}{\hline}\providecommand{\operatornamewithlimits}[1]{\operatorname*{#1}}
\providecommand{\llbracket}{[\![}\providecommand{\rrbracket}{]\!]}\providecommand{\cancel}[1]{#1}
\providecommand{\mathscr}{\mathcal}
% the characters beyond ASCII this document uses
\DeclareUnicodeCharacter{03B1}{\ensuremath{\alpha}}
\DeclareUnicodeCharacter{211D}{\ensuremath{\mathbb{R}}}
\DeclareUnicodeCharacter{2192}{\ensuremath{\rightarrow}}
\DeclareUnicodeCharacter{2194}{\ensuremath{\leftrightarrow}}
\DeclareUnicodeCharacter{21A9}{\ensuremath{\hookleftarrow}}
\DeclareUnicodeCharacter{2208}{\ensuremath{\in}}
\DeclareUnicodeCharacter{2264}{\ensuremath{\leq}}
\DeclareUnicodeCharacter{27E8}{\ensuremath{\langle}}
\DeclareUnicodeCharacter{27E9}{\ensuremath{\rangle}}
% the project's macros ([katex.macros] in blueprint.toml)
\providecommand{\Opens}{}\renewcommand{\Opens}[1]{\mathcal{O}(#1)}
\providecommand{\Top}{}\renewcommand{\Top}{\mathbf{Top}}
\providecommand{\cl}{}\renewcommand{\cl}[1]{\overline{#1}}
\providecommand{\fam}{}\renewcommand{\fam}[1]{\mathcal{#1}}

\begin{document}
\pagestyle{plain}
\begin{center}
{\LARGE\bfseries Compactness through finite intersections\par}
\vspace{0.7em}
{\large The complete blueprint\par}
\vspace{0.5em}
{\small\color{bpmuted}All levels \textperiodcentered{} read along the \texttt{refines} order \textperiodcentered{} appendices included\par}
{\small\color{bpmuted}Snapshot 4c63101a5eb6\par}
\end{center}
\vspace{0.5em}
{\small\tableofcontents}
\bigskip\hrule\bigskip

% object: intro (blueprint/intro.md)
\section*{\texorpdfstring{1\quad Introduction}{1 Introduction}}
\phantomsection\label{bp1}%
\addcontentsline{toc}{section}{\texorpdfstring{1\quad Introduction}{1 Introduction}}
We prove that a space is compact exactly when every family of closed sets with the finite intersection property has a point in common (\hyperref[bp10]{Theorem 2.2.2}), and use it to show that a product of two compact spaces is compact (\hyperref[bp16]{Theorem 3.1}).  The argument is the one in Čech's and Erdős' lecture notes, written so that every step can be checked on its own.

Throughout, \(X\) and \(Y\) are topological spaces and \(\Opens{X}\) is the set of open subsets of \(X\).

% object: spaces (blueprint/spaces/_section.md)
\section*{\texorpdfstring{2\quad Covers and compact spaces}{2 Covers and compact spaces}}
\phantomsection\label{bp2}%
\addcontentsline{toc}{section}{\texorpdfstring{2\quad Covers and compact spaces}{2 Covers and compact spaces}}
The vocabulary of covers, and the characterisation of compactness by closed sets.
% object: covers (blueprint/spaces/covers/_section.md)
\subsection*{\texorpdfstring{2.1\quad Covers}{2.1 Covers}}
\phantomsection\label{bp3}%
\addcontentsline{toc}{subsection}{\texorpdfstring{2.1\quad Covers}{2.1 Covers}}
% text: covers-opening
\phantomsection\label{bp4}%
Two notions recur below, and both are about families of subsets indexed by an arbitrary set \(I\).

% object: def-cover (blueprint/spaces/covers/def-cover.md)
\phantomsection\label{bp5}%
\begin{bpstatement}{Definition~2.1.1}{ (Open cover)}
An \emph{open cover} of \(X\) is a family \(\fam{U} = (U_i)_{i \in I}\) of open sets with

\[
\bigcup_{i \in I} U_i = X .
\]

A \emph{subcover} is a subfamily \((U_j)_{j \in J}\), \(J \subseteq I\), that is still a cover; it is \emph{finite} when \(J\) is.
\end{bpstatement}

% object: def-fip (blueprint/spaces/covers/def-fip.md)
\phantomsection\label{bp6}%
\begin{bpstatement}{Definition~2.1.2}{ (Finite intersection property)}
A family \((F_i)_{i \in I}\) of subsets of \(X\) has the \emph{finite intersection property} when every finite subfamily meets: \[
\begin{aligned}
F_{i_1} \cap \dots \cap F_{i_n} &\neq \emptyset
  &&\text{for all } n \ge 1 \text{ and } i_1, \dots, i_n \in I .
\end{aligned}
\] It is the dual of \hyperref[bp5]{Definition 2.1.1} under complements.
\end{bpstatement}


% object: compact (blueprint/spaces/compact/_section.md)
\subsection*{\texorpdfstring{2.2\quad Compact spaces}{2.2 Compact spaces}}
\phantomsection\label{bp7}%
\addcontentsline{toc}{subsection}{\texorpdfstring{2.2\quad Compact spaces}{2.2 Compact spaces}}
% object: def-compact (blueprint/spaces/compact/def-compact.md)
\phantomsection\label{bp8}%
\begin{bpstatement}{Definition~2.2.1}{ (Compact space)}
\(X\) is \emph{compact} when every open cover of \(X\) (\hyperref[bp5]{Definition 2.1.1}) has a finite subcover.
\end{bpstatement}

% text: compact-between
\phantomsection\label{bp9}%
Passing to complements turns open covers into families of closed sets, and \textquotedbl{}has no finite subcover\textquotedbl{} into \textquotedbl{}has the finite intersection property\textquotedbl{}.  The next theorem records this.

% object: thm-fip (blueprint/spaces/compact/thm-fip.md)
\phantomsection\label{bp10}%
\addcontentsline{toc}{subsubsection}{\texorpdfstring{2.2.2\quad Compactness by closed sets}{2.2.2 Compactness by closed sets}}
\begin{bpstatementit}{Theorem~2.2.2}{ (Compactness by closed sets)}
\(X\) is compact if and only if every family \((F_i)_{i \in I}\) of closed subsets of \(X\) with the finite intersection property has \(\bigcap_{i \in I} F_i \neq \emptyset\).
\end{bpstatementit}
% object: proof/thm-fip (blueprint/spaces/compact/thm-fip.proof.md)
\phantomsection\label{bp11}%
\begin{proof}[{Proof 2.2.2.1}]
{\small\color{bpmuted}Of \hyperref[bp10]{Compactness by closed sets} (Theorem 2.2.2).}\par
Both directions are the contrapositive of each other once the complements are taken: \hyperref[bp12]{Lemma 2.2.2.1.1} exchanges the two conditions, and \hyperref[bp13]{Lemma 2.2.2.1.2} says what a finite subcover becomes.
% step: uses/proof~thm-fip/def-fip
\begin{bpstep}{\textsc{uses} \ensuremath{\to} \hyperref[bp6]{Finite intersection property} (Definition 2.1.2)}\phantomsection\label{bp21}%
\textit{Where the finite intersection property enters.}
Only in \hyperref[bp13]{Lemma 2.2.2.1.2}: the property is the statement that no finite subfamily has empty intersection, word for word.
\end{bpstep}
% object: lem-complements (blueprint/spaces/compact/lem-complements.md)
\phantomsection\label{bp12}%
\begin{bpstatementit}{Lemma~2.2.2.1.1}{ (Complements)}
For a family \((F_i)_{i \in I}\) of closed sets put \(U_i = X \setminus F_i\). Then \[
\bigcap_{i \in I} F_i = \emptyset
\iff
\bigcup_{i \in I} U_i = X ,
\] so \((U_i)\) is an open cover exactly when the \(F_i\) have no common point.
\end{bpstatementit}

% object: lem-finite-subfamily (blueprint/spaces/compact/lem-finite-subfamily.md)
\phantomsection\label{bp13}%
\begin{bpstatementit}{Lemma~2.2.2.1.2}{ (Finite subfamilies)}
For \(J \subseteq I\) finite, \((U_j)_{j \in J}\) covers \(X\) if and only if \(\bigcap_{j \in J} F_j = \emptyset\).  Hence \((U_i)\) has no finite subcover if and only if \((F_i)\) has the finite intersection property.
\end{bpstatementit}

\end{proof}


% object: lem-tube (blueprint/spaces/compact/lem-tube.md)
\phantomsection\label{bp14}%
\begin{bpstatementit}{Lemma~2.2.3}{ (Tube lemma)}
Let \(Y\) be compact, \(x \in X\), and \(W \subseteq X \times Y\) open with \(\{x\} \times Y \subseteq W\).  Then there is an open \(U \ni x\) with \(U \times Y \subseteq W\).
\end{bpstatementit}
% proof: proof/lem-tube
\begin{proof}[{Proof}]\phantomsection\label{bp22}%
Cover \(\{x\} \times Y\) by basic open sets \(U_y \times V_y \subseteq W\).  The \(V_y\) cover \(Y\), so finitely many do, say \(V_{y_1}, \dots, V_{y_n}\), and \(U = U_{y_1} \cap \dots \cap U_{y_n}\) is the open set wanted.
\end{proof}



% object: products (blueprint/products/_section.md)
\section*{\texorpdfstring{3\quad Products}{3 Products}}
\phantomsection\label{bp15}%
\addcontentsline{toc}{section}{\texorpdfstring{3\quad Products}{3 Products}}
Finite products of compact spaces are compact.  The tube lemma (\hyperref[bp14]{Lemma 2.2.3}) is stated in Section \hyperref[bp7]{2.2}; it is the whole content of the argument.
% object: thm-product (blueprint/products/thm-product.md)
\phantomsection\label{bp16}%
\begin{bpstatementit}{Theorem~3.1}{ (Products of two compact spaces)}
If \(X\) and \(Y\) are compact, so is \(X \times Y\) with the product topology.

\begin{proof}[{Proof}]
Let \(\fam{W}\) be an open cover of \(X \times Y\).  For each \(x \in X\) the slice \(\{x\} \times Y \cong Y\) is compact, so finitely many members of \(\fam{W}\) cover it; call their union \(W_x\).  By \hyperref[bp14]{Lemma 2.2.3} there is an open \(U_x \ni x\) with \(U_x \times Y \subseteq W_x\).

\begin{enumerate}
\item The \(U_x\) cover \(X\), so by \hyperref[bp8]{Definition 2.2.1} finitely many do: \(U_{x_1}, \dots, U_{x_m}\).
\item Then \(X \times Y = \bigcup_{k=1}^{m} U_{x_k} \times Y \subseteq \bigcup_k W_{x_k}\), a union of finitely many members of \(\fam{W}\).
\end{enumerate}

Compare the closed-set form, \hyperref[bp10]{Theorem 2.2.2}, which gives the same result by a less direct route.
\end{proof}
\end{bpstatementit}

% object: rem-infinite (blueprint/products/rem-infinite.md)
\phantomsection\label{bp17}%
\begin{bpstatement}{Remark~3.2}{ (Arbitrary products)}
For an arbitrary family \((X_\alpha)_{\alpha \in A}\) the statement is Tychonoff's theorem, which needs the axiom of choice; the argument above only gives finite products, by induction on the number of factors: \(\prod_{k=1}^{n+1} X_k \cong \bigl(\prod_{k=1}^{n} X_k\bigr) \times X_{n+1}\).
\end{bpstatement}

% repeated: lem-tube
\begin{bpnote}\hyperref[bp14]{Lemma 2.2.3} (Tube lemma) belongs here too; it is written out under \hyperref[bp7]{2.2 Compact spaces}, on page~\pageref{bp14}.\end{bpnote}

% object: conventions (blueprint/appendix/_section.md)
\section*{\texorpdfstring{4\quad Conventions}{4 Conventions}}
\phantomsection\label{bp18}%
\addcontentsline{toc}{section}{\texorpdfstring{4\quad Conventions}{4 Conventions}}
Notation used above, collected in one place.
% object: notation (blueprint/appendix/notation.md)
\phantomsection\label{bp19}%
\begin{bpstatement}{Remark~4.1}{ (Symbols)}
\begin{center}\small
\begin{tabular}{llr}
\toprule
Symbol & Meaning & First used \\
\midrule
\(\cl{A}\) & the closure of \(A\) & \hyperref[bp6]{Definition 2.1.2} \\
\(\Opens{X}\) & the open subsets of \(X\) & \hyperref[bp1]{Section 1} \\
\(\fam{U}\), \(\fam{W}\) & families of open sets & \hyperref[bp5]{Definition 2.1.1} \\
\(\Top\) & the category of spaces, \(f\colon X \to Y\) continuous & — \\
\bottomrule
\end{tabular}
\end{center}

Greek letters index families (α ∈ A), and ℝ, ≤, → are used as usual. Identifiers in the Lean sources look like \texttt{TopologicalSpace.\allowbreak{}IsCompact} or \texttt{isCompact\_\allowbreak{}iff\_\allowbreak{}finite\_\allowbreak{}subfamily\_\allowbreak{}closed}; a typical lookup is

\begin{bpverbatim}
example (X : Type*) [TopologicalSpace X] :
    CompactSpace X ↔ IsCompact (Set.univ : Set X) := ⟨fun h => h.1, fun h => ↩
  ⟨h⟩⟩
\end{bpverbatim}

\begin{quote}
Every result here is classical; see any text on point-set topology.
\end{quote}
\end{bpstatement}

% object: figure-cover (blueprint/appendix/figure.md)
\phantomsection\label{bp20}%
\begin{bpstatement}{Remark~4.2}{ (A finite subcover)}
The picture shows three open sets covering a compact interval; two of them already suffice.

\begin{center}
\bpimage{assets/f746f617-cover.png}\par\smallskip{\small Three open intervals covering a closed interval}
\end{center}
\end{bpstatement}



\end{document}
